\documentclass{article}
\usepackage[utf8]{inputenc}
\usepackage[T1]{fontenc}
\usepackage[margin=2cm]{geometry}
\usepackage{hyperref}
\sloppy
\begin{document}
\section*{Historical source}
Encyclopaedia Britannica scan, source row 681. Parse 3.5 extracted the text. Super 3.5 VL corrected it from an enlarged crop of each block. It can still contain recognition errors.
Page headers, footers, and page numbers are removed.

\subsection*{Untitled continuation (E0061)}
\textit{Continues an entry from an earlier page.}

the eye at o: fo that none of the rays that fall above b, or below i, can be reflected to the eye at o; and the dis- tance between b and i is equal to half the length of the object AB.

Hence it appears, that if a man sees his whole image in a plane looking-glass, the part of the glass that reflects his image must be just half as long and half as broad as himself, let him stand at any distance from it whatever; and that his image must appear just as far behind the glass as he is before it. Thus, the man AB (No. 30.) viewing himself in the plane mirror CD, which is just half as long as himself, sees his whole image as at ES, behind the glass, exactly equal to his own size. For, a ray AC, proceeding from his eye at A, and falling perpendicularly upon the surface of the glass at C, is reflected back to his eye, in the same line CA; and the eye of his image will appear at E, in the same line produced to E, beyond the glass. And a ray BD, flowing

Of the DIFFERENT REFRANGIBILITY of LIGHT.

We have hitherto supposed that a particle of light, as it comes from the sun, is the least particle into which light can be separated. But we must now correct this supposition, by showing, that a particle of light, as it comes from the sun, is, properly speaking, a bundle of rays, which may be separated from one another. Therefore, for the future, by a ray of light we must be understood to mean, not that collection of particles which we have hitherto called by this name, but the least particles into which light can be separated.

Rays of light are said to be differently refrangible, when at the same or equal angles of incidence some are more turned out of the way than others.

Rays are said to be differently reflexible, if some are more easily reflected than others.

Light is called homogeneous, when all the rays are equally refrangible: it is called heterogeneous, when some rays are more refrangible than others.

The colours of homogeneous light are called primary or simple colours; and those of heterogeneous light are called secondary or mixed.

The rays of the fun are not all equally refrangible: and those rays, which have a different degree of refrangibility, have likewise a different colour.

If a beam of light SF, (No. 32.) that comes from the fun, passes into a dark room through F a round hole in a window shutter EG; this beam proceeding straight forwards, and falling upon a paper at Y, would there make a round picture of the fun. This picture would be a confused one indeed, if the hole is a large one and there is no lens in the hole. However, as this round spot of light is a picture of the fun, we shall hereafter, notwithstanding its confusion, call it by this name. Now if a glass prism ABC is placed between the hole in the window-shutter and the paper at Y, the rays of this beam, by the refraction which they suffer in the prism, will be bent from their straight course, and instead of going on to as to fall upon the paper at Y, they will be turned upwards, and the picture of the fun produced by them will fall upon a paper MN that is placed above Y. If all the rays

Vol. III. No 87.

object at fg upon a white wall, a sheet or a screen of white paper, provided the object is farther from the lens than its principal focus. To make this picture appear distinct and bright, it must have no other light fall upon it but what comes through the lens kl; and for this reason the whole apparatus is to be placed in a dark room EFGH. The lens kl must be very convex, so that the object de may be very near to it, and yet not be nearer than its principal focus: for by this means, as the object is near to the lens, the picture fg will be at a great distance from it, and consequently the picture will be much bigger than the object. Since the picture is inverted in respect of the object, in order to make the picture appear with the right end upwards, it is necessary that the object de should be placed with the wrong end upwards.

Of the MAGIC LANTERN.

ABCD (No. 31.) is a tin lantern, with a tube nklm fixed in the field of it. This tube consists of two joints, one of which slips into the other: and by drawing this joint out, or pushing it in, the tube may be made longer or shorter. At kl in the end of the moveable joint of the tube a convex lens is fixed, and an object painted with transparent colours upon a piece of thin glass is placed at de somewhere in the immoveable joint of the tube; so that as the tube is lengthened or shortened, the lens will be either at a greater or a less distance from this transparent object. In the field of the lantern there is a very convex lens bbc, which serves to cast a very strong light from the candle within the lantern upon the object de. Now when the rays, which shine through the object de, diverge from the several points as d, e, \&c. in the object, and fall upon the lens kl, they will be made to converge to as many points f, g, \&c. on the other side of the lens, and will paint an inverted picture of the


\section*{Source record}
\begin{description}
\item[Scan ID] \texttt{britannica:309f765b1ecc83e0a7f7762e44b4a927a7555443:train:681}
\item[Image file] \texttt{page-0681-c2be06f2d615.jpg}
\item[Image SHA-256] \texttt{c2be06f2d6151658124b91589cab6163caa72e653bc6e1d62ecc30bbf9a202ee}
\item[Dataset] NationalLibraryOfScotland/encyclopaedia\_britannica\_illustrated at revision \texttt{309f765b1ecc83e0a7f7762e44b4a927a7555443} (CC0-1.0)
\item[Extraction] model \texttt{nvidia/nemotron-3.5-parse}, raw reply \texttt{extract/raw/page-0681-c2be06f2d615.json}
\item[Raw reply SHA-256] \texttt{87c41adb8f0b47c8450489ced90db7710e49b7801397b7f186e6101220ee11e0}
\item[Correction] model \texttt{nvidia/Nemotron-3.5-Super-VL}, block crops, raw replies \texttt{correct/raw/crop\_nothink/page-0681-c2be06f2d615-b*.json}; super\_vl: 14
\item[Running header] O P T I C S. 427
\end{description}
\end{document}
